% archon:covers Run202608192034/Basic.lean
% Historical Archon track only. The canonical Anderson/ proof is audited
% separately by StatementAudit.lean and StrictAxiomAudit.lean.
\chapter{Quasi-completeness and completion criteria}

\begin{definition}[Quasi-complete local ring]
  \label{def:quasi_complete}
  \lean{Run202608192034.TODO.QuasiComplete}
  % SOURCE: references/anderson_2014_fulltext.txt, lines 68--76.
  \textit{Source: D.~D. Anderson, Definition~1,
  references/anderson\_2014\_fulltext.txt, lines 68--76,
  DOI 10.1007/978-1-4939-0925-4\_2.}
  Let \((R,\mathfrak m)\) be a Noetherian local ring.  It is
  \emph{quasi-complete} if, for every decreasing sequence
  \[
    I_1\supseteq I_2\supseteq\cdots
  \]
  of ideals and every integer \(k\geq 1\), there is an index \(s\) for which
  \[
    I_s\subseteq \left(\bigcap_{n\geq 1}I_n\right)+\mathfrak m^k.
  \]
\end{definition}

\begin{proof}
  This is a definition.  The intersection is the infimum of the sequence in
  the lattice of ideals, and the sum and power are the usual ideal operations.
\end{proof}

\begin{definition}[Weakly quasi-complete local ring]
  \label{def:weakly_quasi_complete}
  \lean{Run202608192034.TODO.WeaklyQuasiComplete}
  \uses{def:quasi_complete}
  % SOURCE: references/anderson_2014_fulltext.txt, lines 68--76.
  \textit{Source: D.~D. Anderson, Definition~1,
  references/anderson\_2014\_fulltext.txt, lines 68--76,
  DOI 10.1007/978-1-4939-0925-4\_2.}
  Let \((R,\mathfrak m)\) be a Noetherian local ring.  It is \emph{weakly
  quasi-complete} if every decreasing sequence of ideals with zero
  intersection eventually lies in each prescribed power of \(\mathfrak m\):
  for every \(k\geq 1\), some \(s\) satisfies \(I_s\subseteq\mathfrak m^k\).
\end{definition}

\begin{proof}
  This is a definition, obtained from \cref{def:quasi_complete} by restricting
  to chains whose intersection is the zero ideal.
\end{proof}

\begin{definition}[Generic formal fiber]
  \label{def:generic_formal_fiber}
  \lean{Run202608192034.TODO.genericFormalFiber}
  Let \((A,\mathfrak m)\) be a Noetherian local domain, let
  \(K=\operatorname{Frac}(A)\), and let \(\widehat A\) be the
  \(\mathfrak m\)-adic completion.  The generic formal fiber is
  \[
    \operatorname{Spec}(\widehat A\otimes_A K).
  \]
  Equivalently, its points are the prime ideals \(P\) of \(\widehat A\) with
  \(P\cap A=(0)\).
\end{definition}

\begin{proof}
  Localization at the nonzero elements of \(A\) identifies
  \(\widehat A\otimes_A K\) with that localization of \(\widehat A\).
  Prime ideals of a localization correspond exactly to primes disjoint from
  the localized multiplicative set, which here says \(P\cap A=(0)\).
\end{proof}

\begin{definition}[Analytic irreducibility]
  \label{def:analytically_irreducible}
  \lean{Run202608192034.TODO.AnalyticallyIrreducible}
  A Noetherian local domain \(A\) is \emph{analytically irreducible} when its
  adic completion \(\widehat A\) is an integral domain.
\end{definition}

\begin{proof}
  This is a definition.
\end{proof}

\begin{lemma}[Quotient by a quotient is quotient by a sum]
  \label{lem:quotient_sup_quotient_equiv}
  \lean{Run202608192034.TODO.quotientSupQuotientEquiv}
  If \(I,J\subseteq R\) are ideals, then there is a natural ring equivalence
  \[
    R/(I+J)\cong (R/I)/(J(R/I)).
  \]
\end{lemma}

\begin{proof}
  Use the quotient map \(R\to R/I\) and the induced map from \(R/(I+J)\) to
  \((R/I)/(J(R/I))\).  Its kernel is zero because the preimage of \(J(R/I)\)
  is \(I+J\), using `Ideal.comap_map_of_surjective`; it is surjective because
  \(R\to R/I\) is surjective.  Apply `RingEquiv.ofBijective`.
\end{proof}

\begin{lemma}[The quotient-by-sum equivalence on representatives]
  \label{lem:quotient_sup_quotient_equiv_mk}
  \lean{Run202608192034.TODO.quotientSupQuotientEquiv_mk}
  \uses{lem:quotient_sup_quotient_equiv}
  Under the equivalence \(R/(I+J)\cong (R/I)/(J(R/I))\), the class of
  \(x\in R\) maps to the class of the class of \(x\) in \(R/I\).
\end{lemma}

\begin{proof}
  This follows by unfolding the induced quotient map.
\end{proof}

\begin{lemma}[Quotient characterization of quasi-completeness]
  \label{lem:quotient_weak_characterization}
  \lean{Run202608192034.TODO.quasiComplete_iff_all_quotients_weak}
  \uses{def:quasi_complete,def:weakly_quasi_complete}
  % SOURCE: references/anderson_2014_fulltext.txt, lines 160--167 and 170--262.
  \textit{Source: D.~D. Anderson, discussion before Theorem~4 and
  Theorem~4, references/anderson\_2014\_fulltext.txt, lines 160--167 and
  170--262, DOI 10.1007/978-1-4939-0925-4\_2.}
  A Noetherian local ring \(R\) is quasi-complete if and only if every quotient
  \(R/J\) is weakly quasi-complete.
\end{lemma}

\begin{proof}
  Suppose first that \(R\) is quasi-complete.  A decreasing zero-intersection
  chain in \(R/J\) pulls back to a decreasing chain \((I_n)\) of ideals of
  \(R\) whose intersection is \(J\).  Quasi-completeness gives
  \(I_s\subseteq J+\mathfrak m^k\); passing to the quotient gives the required
  containment in the \(k\)-th power of the maximal ideal of \(R/J\).

  Conversely, given a decreasing chain \((I_n)\) in \(R\), put
  \(J=\bigcap_n I_n\).  The images \(I_n/J\) have zero intersection in \(R/J\).
  Weak quasi-completeness of this quotient gives
  \(I_s/J\subseteq(\mathfrak m/J)^k\).  Powers commute with passage to the
  quotient, so lifting the containment yields
  \(I_s\subseteq J+\mathfrak m^k\), which is quasi-completeness of \(R\).
\end{proof}

\begin{lemma}[Completion-prime criterion]
  \label{lem:farley_completion_criterion}
  \lean{Run202608192034.TODO.weaklyQuasiComplete_iff_completion_primes}
  \uses{def:weakly_quasi_complete,def:generic_formal_fiber}
  % SOURCE: references/farley_2016_fulltext.txt, lines 61--63.
  \textit{Source: J.~D. Farley, Proposition~1,
  references/farley\_2016\_fulltext.txt, lines 61--63,
  DOI 10.4134/BKMS.b140895.}
  A Noetherian local domain \(A\) is weakly quasi-complete if and only if every
  nonzero prime ideal \(P\) of \(\widehat A\) has nonzero contraction to \(A\).
  Equivalently, its generic formal fiber consists only of the zero prime.
\end{lemma}

\begin{proof}
  The proof is the completion criterion underlying Farley's proposition.
  For the forward implication, a nonzero completion prime with zero
  contraction gives, by successive approximation modulo
  \(\mathfrak m^n\widehat A\), a decreasing chain of contraction ideals in
  \(A\) with zero intersection which never enters a fixed power of
  \(\mathfrak m\), contradicting weak quasi-completeness.  For the reverse
  implication, failure of weak quasi-completeness supplies such a chain;
  extend the corresponding proper ideal in the completion to a prime.
  Krull intersection makes its contraction zero, while the failed fixed-power
  containment makes the completion prime nonzero.  The localization
  description in \cref{def:generic_formal_fiber} gives the equivalent
  formulation.
\end{proof}

\begin{definition}[Adic completion prime-contraction condition]
  \label{def:adic_completion_prime_contraction_condition}
  \lean{Run202608192034.TODO.AdicCompletionPrimeContractionCondition}
  \uses{def:generic_formal_fiber}
  For a ring \(A\), an ideal \(\mathfrak m\), and the actual
  \(\mathfrak m\)-adic completion \(\widehat A\), this is the condition that
  every nonzero prime ideal \(P\subseteq\widehat A\) has nonzero contraction
  to \(A\).
\end{definition}

\begin{proof}
  This is a named Lean definition packaging the prime-contraction side of
  \cref{lem:farley_completion_criterion} for the concrete target
  \(\widehat A=\operatorname{AdicCompletion}_{\mathfrak m}(A)\).
\end{proof}

\begin{lemma}[Transporting prime contraction across a completion equivalence]
  \label{lem:adic_completion_prime_contraction_condition_iff_target}
  \lean{Run202608192034.TODO.adicCompletionPrimeContractionCondition_iff_target}
  \uses{def:adic_completion_prime_contraction_condition}
  If \(e:\widehat A\cong T\) and
  \(\iota=e\circ(A\to\widehat A)\), then the nonzero-prime contraction
  condition in \(\widehat A\) is equivalent to the corresponding condition in
  \(T\).
\end{lemma}

\begin{proof}
  Map a completion prime forward along \(e\), and comap a target prime
  backward.  Ring equivalences preserve primality and the bottom ideal, while
  \(\operatorname{comap}_e(\operatorname{map}_e P)=P\) and
  \(\operatorname{map}_e(\operatorname{comap}_e Q)=Q\).  Map compatibility
  then identifies the two contractions to \(A\).
\end{proof}

\begin{definition}[Bad weak-completeness chain]
  \label{def:weakly_quasi_complete_bad_chain}
  \lean{Run202608192034.TODO.WeaklyQuasiCompleteBadChain}
  \uses{def:weakly_quasi_complete}
  A bad chain is a decreasing sequence of ideals \(I_n\) with
  \(\bigcap_n I_n=0\), together with a fixed power \(\mathfrak m^k\), \(k>0\),
  such that no \(I_s\) is contained in \(\mathfrak m^k\).
\end{definition}

\begin{proof}
  This is the explicit negation witness for weak quasi-completeness.  It is
  introduced to make Farley's proof follow the source order: failure of weak
  quasi-completeness is represented by an actual descending chain.
\end{proof}

\begin{definition}[Bad completion prime]
  \label{def:adic_completion_bad_prime}
  \lean{Run202608192034.TODO.AdicCompletionBadPrime}
  \uses{def:adic_completion_prime_contraction_condition}
  A bad completion prime is a nonzero prime
  \(P\subseteq\operatorname{AdicCompletion}_{\mathfrak m}(A)\) whose
  contraction to \(A\) is zero.
\end{definition}

\begin{proof}
  This packages the obstruction to the prime-contraction condition.  In
  Farley's language, it is a nonzero point of the generic formal fiber.
\end{proof}

\begin{definition}[Adic power image in the completion]
  \label{def:adic_completion_power_image}
  \lean{Run202608192034.TODO.adicCompletionPowerImage}
  \uses{def:adic_completion_prime_contraction_condition}
  For each \(n\), this is the image of \(\mathfrak m^n\) in the actual
  adic completion \(\widehat A\).
\end{definition}

\begin{proof}
  This is a Lean definition:
  \(\operatorname{map}(A\to\widehat A)(\mathfrak m^n)\).
\end{proof}

\begin{definition}[Contraction chain attached to a bad prime]
  \label{def:bad_prime_contraction_chain}
  \lean{Run202608192034.TODO.badPrimeContractionChain}
  \uses{def:adic_completion_power_image}
  Given a completion prime \(P\subseteq\widehat A\), define
  \[
    I_n=(P+\mathfrak m^n\widehat A)\cap A.
  \]
  This is Farley's natural descending chain associated to \(P\).
\end{definition}

\begin{proof}
  In Lean this is
  \(\operatorname{comap}(A\to\widehat A)(P\sqcup
  \operatorname{map}(A\to\widehat A)(\mathfrak m^n))\).
\end{proof}

\begin{lemma}[Membership in the contraction chain]
  \label{lem:mem_bad_prime_contraction_chain}
  \lean{Run202608192034.TODO.mem_badPrimeContractionChain}
  \uses{def:bad_prime_contraction_chain}
  An element \(x\in A\) lies in \(I_n\) exactly when its image in
  \(\widehat A\) lies in \(P+\mathfrak m^n\widehat A\).
\end{lemma}

\begin{proof}
  This is checked by unfolding the definition of ideal comap.
\end{proof}

\begin{lemma}[Adic power images are antitone]
  \label{lem:adic_completion_power_image_antitone}
  \lean{Run202608192034.TODO.adicCompletionPowerImage_antitone}
  \uses{def:adic_completion_power_image}
  If \(m\le n\), then the image of \(\mathfrak m^n\) in \(\widehat A\) is
  contained in the image of \(\mathfrak m^m\).
\end{lemma}

\begin{proof}
  This is checked in Lean from Mathlib's
  \texttt{Ideal.pow\_le\_pow\_right} and monotonicity of ideal map.
\end{proof}

\begin{lemma}[The bad-prime contraction chain is descending]
  \label{lem:bad_prime_contraction_chain_antitone}
  \lean{Run202608192034.TODO.badPrimeContractionChain_antitone}
  \uses{def:bad_prime_contraction_chain,lem:adic_completion_power_image_antitone}
  The ideals \(I_n=(P+\mathfrak m^n\widehat A)\cap A\) form a descending
  chain.
\end{lemma}

\begin{proof}
  This is checked in Lean by applying the antitonicity of
  \(\mathfrak m^n\widehat A\), monotonicity of supremum with fixed \(P\), and
  monotonicity of ideal comap.
\end{proof}

\begin{lemma}[The adic power lies in the contraction chain]
  \label{lem:power_le_bad_prime_contraction_chain}
  \lean{Run202608192034.TODO.power_le_badPrimeContractionChain}
  \uses{def:bad_prime_contraction_chain}
  For every \(n\), one has \(\mathfrak m^n\subseteq I_n\).
\end{lemma}

\begin{proof}
  This is checked directly: elements of \(\mathfrak m^n\) map into
  \(\mathfrak m^n\widehat A\), hence into \(P+\mathfrak m^n\widehat A\).
\end{proof}

\begin{lemma}[Prime contraction lies in every chain term]
  \label{lem:comap_le_bad_prime_contraction_chain}
  \lean{Run202608192034.TODO.comap_le_badPrimeContractionChain}
  \uses{def:bad_prime_contraction_chain}
  The contraction \(P\cap A\) is contained in every
  \(I_n=(P+\mathfrak m^n\widehat A)\cap A\).
\end{lemma}

\begin{proof}
  This is checked directly from \(P\subseteq P+\mathfrak m^n\widehat A\) and
  monotonicity of ideal comap.
\end{proof}

\begin{lemma}[Prime contraction lies in the chain intersection]
  \label{lem:comap_le_bad_prime_contraction_chain_iinf}
  \lean{Run202608192034.TODO.comap_le_badPrimeContractionChain_iInf}
  \uses{lem:comap_le_bad_prime_contraction_chain}
  The contraction \(P\cap A\) is contained in \(\bigcap_n I_n\).
\end{lemma}

\begin{proof}
  This is checked in Lean by applying the previous containment for every
  \(n\).
\end{proof}

\begin{lemma}[The first completed power is proper]
  \label{lem:adic_completion_power_image_one_ne_top}
  \lean{Run202608192034.TODO.adicCompletionPowerImage_one_ne_top}
  \uses{def:adic_completion_power_image}
  If \(\mathfrak m\) is the maximal ideal of \(A\), then its image
  \(\mathfrak m\widehat A\) is a proper ideal of \(\widehat A\).
\end{lemma}

\begin{proof}
  This is checked in Lean.  If \(1\in\mathfrak m\widehat A\), then applying
  the finite projection
  \(\widehat A\to A/\mathfrak m\) sends the top ideal to zero.  Hence
  \(1=0\) in \(A/\mathfrak m\), so \(1\in\mathfrak m\), contradicting that
  \(\mathfrak m\) is the local maximal ideal.
\end{proof}

\begin{lemma}[Hausdorff quotient from a Noetherian local completion]
  \label{lem:adic_completion_quotient_hausdorff_of_noetherian_local_completion}
  \lean{Run202608192034.TODO.adicCompletion_quotient_hausdorff_of_noetherian_local_completion}
  \uses{lem:adic_completion_power_image_one_ne_top}
  If \(\widehat A\) is Noetherian local, then every quotient
  \(\widehat A/P\) is Hausdorff for the \(\mathfrak m\widehat A\)-adic
  topology.
\end{lemma}

\begin{proof}
  This is checked in Lean from Mathlib's Krull-intersection theorem for finite
  modules over Noetherian local rings, packaged as
  \texttt{IsHausdorff.of\_isLocalRing}, together with the preceding properness
  of \(\mathfrak m\widehat A\).
\end{proof}

\begin{lemma}[Hausdorff quotient from a Jacobson containment]
  \label{lem:adic_completion_quotient_hausdorff_of_noetherian_jacobson_completion}
  \lean{Run202608192034.TODO.adicCompletion_quotient_hausdorff_of_noetherian_jacobson_completion}
  \uses{def:adic_completion_power_image}
  If \(\widehat A\) is Noetherian and
  \[
    \mathfrak m\widehat A\subseteq \operatorname{Jac}(\widehat A),
  \]
  then every quotient \(\widehat A/P\) is Hausdorff for the
  \(\mathfrak m\widehat A\)-adic topology.
\end{lemma}

\begin{proof}
  This is checked directly from Mathlib's Krull-intersection theorem in the
  form \texttt{IsHausdorff.of\_le\_jacobson}.  It isolates the remaining
  completion-structure work into Noetherianity plus the Jacobson containment.
\end{proof}

\begin{lemma}[Adic completeness gives the Jacobson containment]
  \label{lem:adic_completion_power_image_one_le_jacobson_of_is_adic_complete}
  \lean{Run202608192034.TODO.adicCompletionPowerImage_one_le_jacobson_of_isAdicComplete}
  \uses{def:adic_completion_power_image}
  If \(\widehat A\) is complete for the
  \(\mathfrak m\widehat A\)-adic topology, then
  \(\mathfrak m\widehat A\subseteq\operatorname{Jac}(\widehat A)\).
\end{lemma}

\begin{proof}
  This is checked directly from Mathlib's
  \texttt{IsAdicComplete.le\_jacobson\_bot}.  It shows that the Jacobson
  containment remaining in \cref{lem:adic_completion_quotient_hausdorff} can be
  obtained from `mAhat`-adic completeness.
\end{proof}

\begin{lemma}[Hausdorff quotient of the completion]
  \label{lem:adic_completion_quotient_hausdorff}
  \lean{Run202608192034.TODO.adicCompletion_quotient_hausdorff_source}
  \uses{lem:adic_completion_quotient_hausdorff_of_noetherian_local_completion,lem:adic_completion_quotient_hausdorff_of_noetherian_jacobson_completion,lem:adic_completion_power_image_one_le_jacobson_of_is_adic_complete,lem:adic_completion_quotient_hausdorff_of_noetherian_original_precomplete,lem:adic_completion_quotient_hausdorff_of_noetherian_completion,lem:adic_completion_power_image_one_is_adic_complete_source,lem:adic_completion_is_precomplete_original_source,lem:adic_completion_power_image_one_is_adic_complete_of_original_precomplete,lem:adic_completion_is_local_ring_of_original_precomplete,lem:adic_completion_power_image_one_is_maximal,lem:adic_completion_power_image_one_is_hausdorff}
  % SOURCE: references/farley_2016_fulltext.txt, lines 61--63.
  \textit{Source: J.~D. Farley, Proposition~1,
  references/farley\_2016\_fulltext.txt, lines 61--63,
  DOI 10.4134/BKMS.b140895.}
  The quotient \(\widehat A/P\) is Hausdorff for the topology induced by
  \(\mathfrak m\widehat A\).
\end{lemma}

\begin{proof}
  This standard closed-ideal/Krull-intersection fact is the explicit auxiliary
  boundary `adicCompletion_quotient_hausdorff_external`.  The surrounding Lean
  lemmas prove the result once Noetherianity of the completion is available,
  and they prove all remaining completeness, localness, maximality, and
  Jacobson-containment steps.  Mathlib v4.29 does not provide the general
  theorem that a Noetherian ring's adic completion is Noetherian.  This
  auxiliary boundary is not in the final theorem's dependency closure.
\end{proof}

\begin{lemma}[Neighborhood intersection from Hausdorff quotient]
  \label{lem:adic_completion_neighborhood_iinf_le_prime_of_hausdorff}
  \lean{Run202608192034.TODO.adicCompletion_neighborhood_iInf_le_prime_of_hausdorff}
  \uses{def:adic_completion_power_image,lem:adic_completion_quotient_hausdorff}
  If \(\widehat A/P\) is Hausdorff for the \(\mathfrak m\widehat A\)-adic
  topology, then
  \[
    \bigcap_n(P+\mathfrak m^n\widehat A)\subseteq P.
  \]
\end{lemma}

\begin{proof}
  This is checked in Lean.  An element in every \(P+\mathfrak m^n\widehat A\)
  maps in \(\widehat A/P\) into every power of \(\mathfrak m\widehat A\).  The
  Hausdorff condition forces its image in the quotient to be zero, hence the
  original element lies in \(P\).
\end{proof}

\begin{lemma}[Completion-neighborhood Krull-intersection containment]
  \label{lem:adic_completion_neighborhood_iinf_le_prime}
  \lean{Run202608192034.TODO.adicCompletion_neighborhood_iInf_le_prime_source}
  \uses{lem:adic_completion_quotient_hausdorff,lem:adic_completion_neighborhood_iinf_le_prime_of_hausdorff}
  % SOURCE: references/farley_2016_fulltext.txt, lines 61--63.
  \textit{Source: J.~D. Farley, Proposition~1,
  references/farley\_2016\_fulltext.txt, lines 61--63,
  DOI 10.4134/BKMS.b140895.}
  Under the Noetherian local hypotheses,
  \[
    \bigcap_n(P+\mathfrak m^n\widehat A)\subseteq P.
  \]
\end{lemma}

\begin{proof}
  This is now a checked wrapper around the Hausdorff quotient input and the
  preceding Lean bridge.
\end{proof}

\begin{lemma}[Krull-intersection reverse containment for the bad-prime chain]
  \label{lem:bad_prime_contraction_chain_iinf_le_comap}
  \lean{Run202608192034.TODO.badPrimeContractionChain_iInf_le_comap}
  \uses{def:bad_prime_contraction_chain,lem:adic_completion_neighborhood_iinf_le_prime}
  The intersection of the contracted chain
  \(I_n=(P+\mathfrak m^n\widehat A)\cap A\) is contained in \(P\cap A\).
\end{lemma}

\begin{proof}
  This is now checked in Lean.  If \(x\) lies in all contracted ideals, then
  its image in \(\widehat A\) lies in every \(P+\mathfrak m^n\widehat A\);
  the completion-neighborhood Krull-intersection input puts that image in
  \(P\), hence \(x\in P\cap A\).
\end{proof}

\begin{lemma}[Intersection of the bad-prime contraction chain]
  \label{lem:bad_prime_contraction_chain_iinf_eq_comap}
  \lean{Run202608192034.TODO.badPrimeContractionChain_iInf_eq_comap}
  \uses{lem:bad_prime_contraction_chain_iinf_le_comap,lem:comap_le_bad_prime_contraction_chain_iinf}
  The intersection of \(I_n=(P+\mathfrak m^n\widehat A)\cap A\) is exactly
  \(P\cap A\).
\end{lemma}

\begin{proof}
  This is now a checked Lean wrapper: one inclusion is formal lattice
  reasoning, and the other is the remaining Krull-intersection source input.
\end{proof}

\begin{lemma}[Zero intersection for a bad-prime contraction chain]
  \label{lem:bad_prime_contraction_chain_iinf_eq_bot}
  \lean{Run202608192034.TODO.badPrimeContractionChain_iInf_eq_bot}
  \uses{lem:bad_prime_contraction_chain_iinf_eq_comap,def:adic_completion_bad_prime}
  If \(P\cap A=0\), then the associated contraction chain has zero
  intersection.
\end{lemma}

\begin{proof}
  This is checked in Lean from
  \cref{lem:bad_prime_contraction_chain_iinf_eq_comap} and the zero
  contraction hypothesis.
\end{proof}

\begin{lemma}[Nonzero completion elements leave some power image]
  \label{lem:adic_completion_nonzero_not_mem_power_image}
  \lean{Run202608192034.TODO.adicCompletion_nonzero_not_mem_powerImage_source}
  \uses{def:adic_completion_power_image}
  If \(0\ne y\in\widehat A\), then there is some \(k>0\) such that
  \(y\notin \mathfrak m^k\widehat A\).
\end{lemma}

\begin{proof}
  This is now checked in Lean.  The proof uses Mathlib's Hausdorff instance
  for \(\operatorname{AdicCompletion}_{\mathfrak m}(A)\), together with
  \texttt{Ideal.smul\_top\_eq\_map} to identify the Hausdorff filtration
  \(\mathfrak m^n\cdot\widehat A\) with the ideal image of
  \(\mathfrak m^n\) in \(\widehat A\).
\end{proof}

\begin{definition}[Finite quotient linear map]
  \label{def:quotient_power_linear_map}
  \lean{Run202608192034.TODO.quotientPowerLinearMap}
  For each \(t\), let
  \[
    A\to A/\mathfrak m^t
  \]
  be the quotient map, regarded as an \(A\)-linear map.
\end{definition}

\begin{proof}
  This is the linear map underlying the usual quotient algebra map
  \(A\to A/\mathfrak m^t\).
\end{proof}

\begin{lemma}[Finite quotient map is surjective]
  \label{lem:quotient_power_linear_map_surjective}
  \lean{Run202608192034.TODO.quotientPowerLinearMap_surjective}
  \uses{def:quotient_power_linear_map}
  The map \(A\to A/\mathfrak m^t\) is surjective.
\end{lemma}

\begin{proof}
  This is checked directly from the quotient representative theorem
  `Ideal.Quotient.mk\_surjective`.
\end{proof}

\begin{lemma}[Exactness after completing the finite quotient sequence]
  \label{lem:adic_completion_power_quotient_exact}
  \lean{Run202608192034.TODO.adicCompletion_power_quotient_exact}
  \uses{def:quotient_power_linear_map,lem:quotient_power_linear_map_surjective}
  Completing the exact sequence
  \[
    \mathfrak m^t\hookrightarrow A\to A/\mathfrak m^t
  \]
  gives an exact pair
  \[
    \widehat{\mathfrak m^t}\to\widehat A
      \to \widehat{A/\mathfrak m^t}.
  \]
\end{lemma}

\begin{proof}
  This is now checked using Mathlib's Noetherian exactness theorem
  `AdicCompletion.map\_exact`, applied to the submodule inclusion
  \(\mathfrak m^t\hookrightarrow A\) and the quotient map.
\end{proof}

\begin{lemma}[Tensor image of the completed power lies in the extended power]
  \label{lem:adic_completion_power_tensor_image_le_power_image}
  \lean{Run202608192034.TODO.adicCompletion_powerTensorImage_le_powerImage}
  \uses{def:adic_completion_power_image}
  The tensor presentation map
  \[
    \widehat A\otimes_A\mathfrak m^t\to\widehat A
  \]
  has image contained in \(\mathfrak m^t\widehat A\).
\end{lemma}

\begin{proof}
  This is checked by tensor induction.  On pure tensors, the image is
  \(r\cdot a\) with \(a\in\mathfrak m^t\), hence lies in the extended ideal;
  the zero and additive cases follow from ideal laws.
\end{proof}

\begin{lemma}[Completed power image lies in the extended power]
  \label{lem:adic_completion_power_submodule_range_le_power_image}
  \lean{Run202608192034.TODO.adicCompletion_powerSubmoduleRange_le_powerImage}
  \uses{lem:adic_completion_power_tensor_image_le_power_image}
  The range of the completed inclusion
  \(\widehat{\mathfrak m^t}\to\widehat A\) is contained in
  \(\mathfrak m^t\widehat A\).
\end{lemma}

\begin{proof}
  This is checked using the surjective tensor presentation
  `AdicCompletion.ofTensorProduct\_surjective\_of\_finite` for the finite
  \(A\)-module \(\mathfrak m^t\), together with naturality of
  `AdicCompletion.ofTensorProduct`.
\end{proof}

\begin{lemma}[Finite projection zero implies completed quotient map zero]
  \label{lem:adic_completion_map_quotient_power_eq_zero_of_eval_zero}
  \lean{Run202608192034.TODO.adicCompletion_map_quotientPower_eq_zero_of_evalₐ_eq_zero}
  \uses{def:quotient_power_linear_map}
  If \(x\in\widehat A\) maps to zero in \(A/\mathfrak m^t\), then its image in
  the completion of \(A/\mathfrak m^t\) is zero.
\end{lemma}

\begin{proof}
  This is checked coordinatewise in the inverse-limit definition of
  `AdicCompletion`.  For levels \(n\le t\), compatibility with transition maps
  propagates the zero \(t\)-th coordinate downward.  For \(t\le n\), a
  representative whose \(t\)-th image is zero already lies in
  \(\mathfrak m^t\), hence maps to zero in \(A/\mathfrak m^t\).
\end{proof}

\begin{lemma}[Kernel of a finite adic projection]
  \label{lem:adic_completion_eval_kernel_le_power_image}
  \lean{Run202608192034.TODO.adicCompletion_evalₐ_kernel_le_powerImage_source}
  \uses{lem:adic_completion_power_quotient_exact,lem:adic_completion_power_submodule_range_le_power_image,lem:adic_completion_map_quotient_power_eq_zero_of_eval_zero}
  For every finite level \(t\), the kernel of the canonical projection
  \[
    \widehat A\to A/\mathfrak m^t
  \]
  is contained in the extended power image
  \(\mathfrak m^t\widehat A\).
\end{lemma}

\begin{proof}
  This is now checked in Lean.  If \(x\) maps to zero at finite level \(t\),
  the preceding coordinate lemma shows that \(x\) maps to zero in
  \(\widehat{A/\mathfrak m^t}\).  Exactness of the completed quotient sequence
  puts \(x\) in the range of
  \(\widehat{\mathfrak m^t}\to\widehat A\), and the tensor-presentation lemma
  identifies that range with a subideal of \(\mathfrak m^t\widehat A\).
\end{proof}

\begin{lemma}[Completed power image is the finite projection kernel]
  \label{lem:adic_completion_power_image_eq_eval_kernel}
  \lean{Run202608192034.TODO.adicCompletionPowerImage_eq_evalₐ_kernel}
  \uses{def:adic_completion_power_image,lem:adic_completion_eval_kernel_le_power_image}
  For every \(t\), the extended ideal
  \(\mathfrak m^t\widehat A\) is exactly the kernel of the finite-level
  projection \(\widehat A\to A/\mathfrak m^t\).
\end{lemma}

\begin{proof}
  The forward containment is checked directly from the formula for
  \texttt{AdicCompletion.eval\_a} on elements coming from \(A\).  The reverse
  containment is the completed exactness argument in
  \cref{lem:adic_completion_eval_kernel_le_power_image}.
\end{proof}

\begin{lemma}[First completed power is the first projection kernel]
  \label{lem:adic_completion_power_image_one_eq_eval_kernel}
  \lean{Run202608192034.TODO.adicCompletionPowerImage_one_eq_evalₐ_kernel}
  \uses{lem:adic_completion_power_image_eq_eval_kernel}
  The extended ideal \(\mathfrak m\widehat A\) is the kernel of the first
  finite-level projection \(\widehat A\to A/\mathfrak m\).
\end{lemma}

\begin{proof}
  This is the preceding all-level kernel equality specialized to \(t=1\).
\end{proof}

\begin{lemma}[Powers of the completed first image]
  \label{lem:adic_completion_power_image_one_pow}
  \lean{Run202608192034.TODO.adicCompletionPowerImage_one_pow}
  \uses{def:adic_completion_power_image}
  For every \(t\), \((\mathfrak m\widehat A)^t=\mathfrak m^t\widehat A\).
\end{lemma}

\begin{proof}
  This is checked by unfolding the extended ideal and using that ideal maps
  commute with powers.
\end{proof}

\begin{lemma}[Ideal action on the top submodule]
  \label{lem:ideal_smul_top_eq_self}
  \lean{Run202608192034.TODO.ideal_smul_top_eq_self}
  For any commutative ring \(B\) and ideal \(J\), the submodule
  \(J\cdot B\) is just \(J\) itself.
\end{lemma}

\begin{proof}
  One containment is ideal absorption under multiplication; the other writes
  \(x\in J\) as \(x\cdot 1\).
\end{proof}

\begin{lemma}[Completed and original finite neighborhoods agree]
  \label{lem:adic_completion_power_image_one_pow_smul_top_restrict_scalars}
  \lean{Run202608192034.TODO.adicCompletionPowerImage_one_pow_smul_top_restrictScalars}
  \uses{lem:adic_completion_power_image_one_pow,lem:ideal_smul_top_eq_self}
  After restricting scalars from \(\widehat A\) to \(A\), the neighborhood
  \((\mathfrak m\widehat A)^t\widehat A\) is the same submodule as
  \(\mathfrak m^t\widehat A\).
\end{lemma}

\begin{proof}
  This is a checked bookkeeping lemma: combine
  \((\mathfrak m\widehat A)^t=\mathfrak m^t\widehat A\) with Mathlib's
  identification of \(I\)-multiples of the top module with the mapped ideal.
\end{proof}

\begin{definition}[Finite-level quotient comparison]
  \label{def:adic_completion_quotient_power_image_equiv}
  \lean{Run202608192034.TODO.adicCompletionQuotientPowerImageEquiv}
  \uses{lem:adic_completion_power_image_eq_eval_kernel}
  The kernel equality gives a checked ring equivalence
  \[
    \widehat A/\mathfrak m^t\widehat A \simeq A/\mathfrak m^t.
  \]
\end{definition}

\begin{proof}
  This equivalence is checked in Lean using quotient-by-equal-ideal
  equivalences and the first isomorphism theorem for the surjective finite
  projection \(\widehat A\to A/\mathfrak m^t\).
\end{proof}

\begin{definition}[Finite-level quotient comparison for powers of \(\mathfrak m\widehat A\)]
  \label{def:adic_completion_quotient_power_equiv}
  \lean{Run202608192034.TODO.adicCompletionQuotientPowerEquiv}
  \uses{lem:adic_completion_power_image_one_pow,def:adic_completion_quotient_power_image_equiv}
  Combining the previous comparison with
  \((\mathfrak m\widehat A)^t=\mathfrak m^t\widehat A\) gives
  \[
    \widehat A/(\mathfrak m\widehat A)^t \simeq A/\mathfrak m^t.
  \]
\end{definition}

\begin{proof}
  This is the previous finite-level quotient comparison after rewriting
  \((\mathfrak m\widehat A)^t\) as \(\mathfrak m^t\widehat A\).
\end{proof}

\begin{lemma}[The first completed power is maximal]
  \label{lem:adic_completion_power_image_one_is_maximal}
  \lean{Run202608192034.TODO.adicCompletionPowerImage_one_isMaximal}
  \uses{lem:adic_completion_power_image_one_eq_eval_kernel}
  If \(\mathfrak m\) is the maximal ideal of \(A\), then
  \(\mathfrak m\widehat A\) is a maximal ideal of \(\widehat A\).
\end{lemma}

\begin{proof}
  This is checked in Lean.  The previous lemma identifies
  \(\mathfrak m\widehat A\) with the kernel of the surjective map
  \(\widehat A\to A/\mathfrak m\), and the target is a field because
  \(\mathfrak m\) is maximal.
\end{proof}

\begin{lemma}[The completed ring is Hausdorff for the completed maximal ideal]
  \label{lem:adic_completion_power_image_one_is_hausdorff}
  \lean{Run202608192034.TODO.adicCompletionPowerImage_one_isHausdorff}
  \uses{def:adic_completion_power_image}
  The completion \(\widehat A\) is Hausdorff for the
  \(\mathfrak m\widehat A\)-adic topology.
\end{lemma}

\begin{proof}
  This is checked in Lean by transporting Mathlib's Hausdorff instance for
  \(\operatorname{AdicCompletion}_{\mathfrak m}(A)\) along
  \(\mathfrak m\widehat A=\operatorname{map}(A\to\widehat A)(\mathfrak m)\).
\end{proof}

\begin{lemma}[Finite-level evaluation kills the matching adic neighborhood]
  \label{lem:adic_completion_eval_eq_of_smod_eq}
  \lean{Run202608192034.TODO.adicCompletion_eval_eq_of_smodEq}
  The finite projection
  \(\widehat A\to A/\mathfrak m^n\) is constant on cosets modulo
  \(\mathfrak m^n\widehat A\), viewed as an \(A\)-submodule of
  \(\widehat A\).
\end{lemma}

\begin{proof}
  The proof is by induction on membership in
  \(\mathfrak m^n\cdot\widehat A\).  A scalar from \(\mathfrak m^n\) kills
  every element after projection to \(A/\mathfrak m^n\), and the kernel is
  closed under addition.
\end{proof}

\begin{lemma}[The adic completion is precomplete for the original ideal]
  \label{lem:adic_completion_is_precomplete_original_source}
  \lean{Run202608192034.TODO.adicCompletion_isPrecomplete_original_source}
  \uses{lem:adic_completion_eval_eq_of_smod_eq,lem:adic_completion_eval_kernel_le_power_image}
  If \(A\) is Noetherian, then \(\widehat A\), as an \(A\)-module, is
  precomplete for the original \(\mathfrak m\)-adic topology.
\end{lemma}

\begin{proof}
  Given a Cauchy sequence \(f_n\) in \(\widehat A\), define the candidate
  limit \(L\) by taking its \(n\)-th finite coordinate to be the image of
  \(f_n\) in \(A/\mathfrak m^n\).  The previous lemma proves these coordinates
  are compatible.  The difference \(f_n-L\) has zero \(n\)-th coordinate, and
  the checked finite-level kernel containment places it in
  \(\mathfrak m^n\widehat A\).
\end{proof}

\begin{lemma}[Precompleteness transfers to the completed maximal ideal]
  \label{lem:adic_completion_power_image_one_is_precomplete_of_original_precomplete}
  \lean{Run202608192034.TODO.adicCompletionPowerImage_one_isPrecomplete_of_original_precomplete}
  \uses{lem:adic_completion_power_image_one_pow_smul_top_restrict_scalars}
  If \(\widehat A\) is precomplete for the original \(\mathfrak m\)-adic
  module structure, then it is precomplete for the
  \(\mathfrak m\widehat A\)-adic structure.
\end{lemma}

\begin{proof}
  This is checked by rewriting the Cauchy and convergence neighborhoods via
  \cref{lem:adic_completion_power_image_one_pow_smul_top_restrict_scalars}.
\end{proof}

\begin{lemma}[The completed maximal ideal is precomplete]
  \label{lem:adic_completion_power_image_one_is_precomplete_source}
  \lean{Run202608192034.TODO.adicCompletionPowerImage_one_isPrecomplete_source}
  \uses{lem:adic_completion_is_precomplete_original_source,lem:adic_completion_power_image_one_is_precomplete_of_original_precomplete}
  If \(A\) is Noetherian, then \(\widehat A\) is precomplete for the
  \(\mathfrak m\widehat A\)-adic topology.
\end{lemma}

\begin{proof}
  Use the original \(\mathfrak m\)-adic precompleteness of \(\widehat A\) and
  the checked transfer lemma comparing original and completed neighborhoods.
\end{proof}

\begin{lemma}[Adic completeness transfers to the completed maximal ideal]
  \label{lem:adic_completion_power_image_one_is_adic_complete_of_original_precomplete}
  \lean{Run202608192034.TODO.adicCompletionPowerImage_one_isAdicComplete_of_original_precomplete}
  \uses{lem:adic_completion_power_image_one_is_hausdorff,lem:adic_completion_power_image_one_is_precomplete_of_original_precomplete}
  Under the same original precompleteness hypothesis, \(\widehat A\) is
  complete for the \(\mathfrak m\widehat A\)-adic topology.
\end{lemma}

\begin{proof}
  This combines the already checked Hausdorffness with the preceding
  precompleteness-transfer lemma.
\end{proof}

\begin{lemma}[The completed maximal-ideal topology is complete]
  \label{lem:adic_completion_power_image_one_is_adic_complete_source}
  \lean{Run202608192034.TODO.adicCompletionPowerImage_one_isAdicComplete_source}
  \uses{lem:adic_completion_power_image_one_is_hausdorff,lem:adic_completion_is_precomplete_original_source,lem:adic_completion_power_image_one_is_precomplete_of_original_precomplete}
  If \(A\) is Noetherian, then \(\widehat A\) is complete for the
  \(\mathfrak m\widehat A\)-adic topology.
\end{lemma}

\begin{proof}
  Use the newly checked original precompleteness of \(\widehat A\), transfer
  it to the completed maximal-ideal filtration, and combine with the checked
  Hausdorffness of \(\widehat A\).
\end{proof}

\begin{lemma}[The completion is local under original precompleteness]
  \label{lem:adic_completion_is_local_ring_of_original_precomplete}
  \lean{Run202608192034.TODO.adicCompletion_isLocalRing_of_original_precomplete}
  \uses{lem:adic_completion_power_image_one_is_maximal,lem:adic_completion_power_image_one_is_adic_complete_of_original_precomplete}
  If \(\mathfrak m\) is the maximal ideal of \(A\) and \(\widehat A\) is
  precomplete for the original \(\mathfrak m\)-adic structure, then
  \(\widehat A\) is a local ring.
\end{lemma}

\begin{proof}
  This is checked from Mathlib's theorem that a ring complete with respect to
  a maximal ideal is local.
\end{proof}

\begin{lemma}[The completion is local]
  \label{lem:adic_completion_is_local_ring_source}
  \lean{Run202608192034.TODO.adicCompletion_isLocalRing_source}
  \uses{lem:adic_completion_power_image_one_is_maximal,lem:adic_completion_power_image_one_is_adic_complete_source}
  If \(A\) is Noetherian local and \(\mathfrak m\) is its maximal ideal, then
  \(\widehat A\) is local.
\end{lemma}

\begin{proof}
  Apply Mathlib's theorem that a ring complete with respect to a maximal ideal
  is local, using the checked maximality of \(\mathfrak m\widehat A\) and the
  checked completed-topology completeness.
\end{proof}

\begin{lemma}[Hausdorff quotient from Noetherianity and original precompleteness]
  \label{lem:adic_completion_quotient_hausdorff_of_noetherian_original_precomplete}
  \lean{Run202608192034.TODO.adicCompletion_quotient_hausdorff_of_noetherian_original_precomplete}
  \uses{lem:adic_completion_quotient_hausdorff_of_noetherian_jacobson_completion,lem:adic_completion_power_image_one_le_jacobson_of_is_adic_complete,lem:adic_completion_power_image_one_is_adic_complete_of_original_precomplete}
  If \(\widehat A\) is Noetherian and precomplete for the original
  \(\mathfrak m\)-adic structure, then every quotient \(\widehat A/P\) is
  Hausdorff for the \(\mathfrak m\widehat A\)-adic topology.
\end{lemma}

\begin{proof}
  The preceding completeness-transfer lemma gives `mAhat`-adic completeness,
  Mathlib gives the Jacobson containment, and
  \cref{lem:adic_completion_quotient_hausdorff_of_noetherian_jacobson_completion}
  applies Krull intersection to the quotient.
\end{proof}

\begin{lemma}[Hausdorff quotient from Noetherianity of the completion]
  \label{lem:adic_completion_quotient_hausdorff_of_noetherian_completion}
  \lean{Run202608192034.TODO.adicCompletion_quotient_hausdorff_of_noetherian_completion}
  \uses{lem:adic_completion_quotient_hausdorff_of_noetherian_jacobson_completion,lem:adic_completion_power_image_one_le_jacobson_of_is_adic_complete,lem:adic_completion_power_image_one_is_adic_complete_source}
  If \(A\) is Noetherian and \(\widehat A\) is Noetherian, then every
  quotient \(\widehat A/P\) is Hausdorff for the
  \(\mathfrak m\widehat A\)-adic topology.
\end{lemma}

\begin{proof}
  The completed-topology completeness gives the Jacobson containment
  \(\mathfrak m\widehat A\subseteq\operatorname{Jac}(\widehat A)\).  The
  Noetherian-Jacobson quotient-Hausdorff lemma then applies.
\end{proof}

\begin{lemma}[Finite-level approximation in the completion]
  \label{lem:adic_completion_exists_approx_mod_power_image}
  \lean{Run202608192034.TODO.adicCompletion_exists_approx_mod_powerImage_source}
  \uses{lem:adic_completion_eval_kernel_le_power_image}
  For every \(y\in\widehat A\) and every finite level \(t\), there is
  \(a\in A\) such that
  \(y-a\in \mathfrak m^t\widehat A\).
\end{lemma}

\begin{proof}
  This is now a checked Lean wrapper.  The finite-level projection
  \(\widehat A\to A/\mathfrak m^t\) is surjective by construction of
  \(\widehat A\); choose a representative \(a\in A\) of the image of
  \(y\).  The difference \(y-a\) lies in the projection kernel, and
  \cref{lem:adic_completion_eval_kernel_le_power_image} puts this kernel in
  \(\mathfrak m^t\widehat A\).
\end{proof}

\begin{lemma}[An approximable nonzero completion element forces fixed-power avoidance]
  \label{lem:bad_prime_contraction_chain_avoids_fixed_power_of_element}
  \lean{Run202608192034.TODO.badPrimeContractionChain_avoids_fixed_power_of_element}
  \uses{def:bad_prime_contraction_chain,lem:adic_completion_power_image_antitone}
  Suppose \(y\in P\), \(y\notin\mathfrak m^k\widehat A\), and \(y\) can be
  approximated modulo every \(\mathfrak m^t\widehat A\) by an element of
  \(A\).  Then the chain
  \(I_s=(P+\mathfrak m^s\widehat A)\cap A\) is not eventually contained in
  \(\mathfrak m^k\); in fact no \(I_s\subseteq\mathfrak m^k\).
\end{lemma}

\begin{proof}
  This is now checked in Lean.  Given \(s\), approximate \(y\) modulo
  \(\mathfrak m^{\max(s,k)}\widehat A\) by \(a\in A\).  Then
  \(a\in I_s\).  If \(a\in\mathfrak m^k\), the approximation error and
  \(a\) both lie in \(\mathfrak m^k\widehat A\), forcing
  \(y\in\mathfrak m^k\widehat A\), contradiction.
\end{proof}

\begin{lemma}[Fixed-power avoidance for the bad-prime chain]
  \label{lem:bad_prime_contraction_chain_avoids_fixed_power}
  \lean{Run202608192034.TODO.badPrimeContractionChain_avoids_fixed_power_source}
  \uses{def:bad_prime_contraction_chain,def:adic_completion_bad_prime,lem:adic_completion_nonzero_not_mem_power_image,lem:adic_completion_exists_approx_mod_power_image,lem:bad_prime_contraction_chain_avoids_fixed_power_of_element}
  % SOURCE: references/farley_2016_fulltext.txt, lines 61--63.
  \textit{Source: J.~D. Farley, Proposition~1,
  references/farley\_2016\_fulltext.txt, lines 61--63,
  DOI 10.4134/BKMS.b140895.}
  If \(P\) is a nonzero completion prime with zero contraction, then the
  associated chain avoids some fixed power of \(\mathfrak m\).
\end{lemma}

\begin{proof}
  This is now a checked Lean wrapper.  Lean chooses a nonzero
  \(y\in P\), applies
  \cref{lem:adic_completion_nonzero_not_mem_power_image} to find the
  obstructing power \(\mathfrak m^k\), and then applies
  \cref{lem:bad_prime_contraction_chain_avoids_fixed_power_of_element}.  The
  remaining source input for this step is exactly the finite-level
  approximation lemma
  \cref{lem:adic_completion_exists_approx_mod_power_image}.
\end{proof}

\begin{lemma}[Failure of weak quasi-completeness is a bad chain]
  \label{lem:not_weakly_quasi_complete_iff_bad_chain}
  \lean{Run202608192034.TODO.not_weaklyQuasiComplete_iff_badChain}
  \uses{def:weakly_quasi_complete,def:weakly_quasi_complete_bad_chain}
  For any ring \(A\) and ideal \(\mathfrak m\),
  \(\neg\operatorname{WeaklyQuasiComplete}(A,\mathfrak m)\) is equivalent to
  the existence of a bad weak-completeness chain.
\end{lemma}

\begin{proof}
  This is a checked Lean proof by unfolding the definition of weak
  quasi-completeness and pushing negations through the quantifiers.
\end{proof}

\begin{lemma}[Failure of prime contraction is a bad completion prime]
  \label{lem:not_adic_completion_prime_contraction_condition_iff_bad_prime}
  \lean{Run202608192034.TODO.not_adicCompletionPrimeContractionCondition_iff_badPrime}
  \uses{def:adic_completion_prime_contraction_condition,def:adic_completion_bad_prime}
  The negation of the adic-completion prime-contraction condition is
  equivalent to the existence of a bad completion prime.
\end{lemma}

\begin{proof}
  This is a checked Lean proof by unfolding
  \(\operatorname{AdicCompletionPrimeContractionCondition}\) and pushing
  negations through the universal quantifiers.
\end{proof}

\begin{lemma}[A bad completion prime gives a bad chain]
  \label{lem:adic_completion_bad_prime_to_bad_chain}
  \lean{Run202608192034.TODO.adicCompletion_badPrime_to_badChain_source}
  \uses{def:adic_completion_bad_prime,def:weakly_quasi_complete_bad_chain,lem:bad_prime_contraction_chain_antitone,lem:bad_prime_contraction_chain_iinf_eq_bot,lem:bad_prime_contraction_chain_avoids_fixed_power}
  % SOURCE: references/farley_2016_fulltext.txt, lines 61--63.
  \textit{Source: J.~D. Farley, Proposition~1,
  references/farley\_2016\_fulltext.txt, lines 61--63,
  DOI 10.4134/BKMS.b140895.}
  For a Noetherian local domain \(A\), a nonzero completion prime with zero
  contraction produces a bad descending chain of ideals in \(A\).
\end{lemma}

\begin{proof}
  This is now a checked Lean wrapper around the explicit chain
  \(I_n=(P+\mathfrak m^n\widehat A)\cap A\).  Lean checks the descending-chain
  property directly; the remaining source inputs are the
  completion-neighborhood Krull-intersection containment and the finite-level
  approximation/kernel-identification statement for adic completion.
\end{proof}

\begin{lemma}[A bad chain gives a bad completion prime]
  \label{lem:weakly_quasi_complete_bad_chain_to_adic_completion_bad_prime}
  \lean{Run202608192034.TODO.weaklyQuasiComplete_badChain_to_adicCompletion_badPrime_source}
  \uses{def:weakly_quasi_complete_bad_chain,def:adic_completion_bad_prime}
  % SOURCE: references/farley_2016_fulltext.txt, lines 61--63.
  \textit{Source: J.~D. Farley, Proposition~1,
  references/farley\_2016\_fulltext.txt, lines 61--63,
  DOI 10.4134/BKMS.b140895.}
  For a Noetherian local domain \(A\), a bad descending chain of ideals
  produces a nonzero prime of the adic completion whose contraction to \(A\)
  is zero.
\end{lemma}

\begin{proof}
  Lean derives this implication from the two checked negation equivalences and
  the explicit source boundary `anderson_corollary2_part1_external`, namely
  Anderson, Corollary~2(1), which Farley quotes as Proposition~1.  The reverse
  bad-prime-to-chain construction is formalized independently above.
\end{proof}

\begin{lemma}[Weak quasi-completeness implies prime contraction]
  \label{lem:weakly_quasi_complete_to_adic_completion_prime_contraction}
  \lean{Run202608192034.TODO.weaklyQuasiComplete_to_adicCompletion_primeContraction_source}
  \uses{def:weakly_quasi_complete,def:adic_completion_prime_contraction_condition,lem:not_weakly_quasi_complete_iff_bad_chain,lem:not_adic_completion_prime_contraction_condition_iff_bad_prime,lem:adic_completion_bad_prime_to_bad_chain}
  Under Farley's hypotheses, weak quasi-completeness implies the
  adic-completion prime-contraction condition.
\end{lemma}

\begin{proof}
  This part is checked in Lean from the preceding source construction.  If
  prime contraction failed, there would be a bad completion prime; Farley's
  first construction turns it into a bad chain, contradicting weak
  quasi-completeness.
\end{proof}

\begin{lemma}[Prime contraction implies weak quasi-completeness]
  \label{lem:adic_completion_prime_contraction_to_weakly_quasi_complete}
  \lean{Run202608192034.TODO.adicCompletion_primeContraction_to_weaklyQuasiComplete_source}
  \uses{def:weakly_quasi_complete,def:adic_completion_prime_contraction_condition,lem:not_weakly_quasi_complete_iff_bad_chain,lem:not_adic_completion_prime_contraction_condition_iff_bad_prime,lem:weakly_quasi_complete_bad_chain_to_adic_completion_bad_prime}
  Under Farley's hypotheses, the adic-completion prime-contraction condition
  implies weak quasi-completeness.
\end{lemma}

\begin{proof}
  This part is checked in Lean from the preceding source construction.  If
  weak quasi-completeness failed, there would be a bad chain; Farley's second
  construction turns it into a bad completion prime, contradicting prime
  contraction.
\end{proof}

\begin{lemma}[Farley's criterion for the actual adic completion]
  \label{lem:weakly_quasi_complete_iff_adic_completion_prime_contraction}
  \lean{Run202608192034.TODO.weaklyQuasiComplete_iff_adicCompletion_primeContraction_source}
  \uses{lem:weakly_quasi_complete_to_adic_completion_prime_contraction,lem:adic_completion_prime_contraction_to_weakly_quasi_complete}
  % SOURCE: references/farley_2016_fulltext.txt, lines 61--63.
  \textit{Source: J.~D. Farley, Proposition~1,
  references/farley\_2016\_fulltext.txt, lines 61--63,
  DOI 10.4134/BKMS.b140895.}
  For a Noetherian local domain \(A\), with the adic ideal identified as the
  maximal ideal, weak quasi-completeness is equivalent to the
  prime-contraction condition for its actual adic completion.
\end{lemma}

\begin{proof}
  This is a checked Lean wrapper.  One implication uses the fully expanded
  bad-prime contraction-chain construction; the other uses the cited
  Corollary~2(1) boundary through the checked negation lemmas.
\end{proof}

\begin{lemma}[Dimension-one prime-contraction/analytic-irreducibility bridge]
  \label{lem:dimension_one_adic_completion_prime_contraction_iff_analytic_irreducible}
  \lean{Run202608192034.TODO.dimensionOne_adicCompletion_primeContraction_iff_analyticIrreducible_source}
  \uses{def:adic_completion_prime_contraction_condition,def:analytically_irreducible}
  % SOURCE: references/anderson_2014_fulltext.txt, lines 307--339.
  \textit{Source: D.~D. Anderson, Corollary~2,
  references/anderson\_2014\_fulltext.txt, lines 307--339,
  DOI 10.1007/978-1-4939-0925-4\_2.}
  In a one-dimensional Noetherian local domain, for the maximal-adic
  completion, the adic-completion prime-contraction condition is equivalent to
  the adic completion being a domain.
\end{lemma}

\begin{proof}
  Lean composes the explicit Anderson Corollary~2(1) and Corollary~2(3)
  boundaries.  The displayed minimal-prime/going-down argument records the
  mathematical proof that would be needed to internalize this bridge.
\end{proof}

\begin{lemma}[Dimension-one analytic criterion for adic completions]
  \label{lem:dimension_one_weak_criterion_adic_completion_source}
  \lean{Run202608192034.TODO.dimensionOne_weaklyQuasiComplete_iff_adicCompletion_source}
  \uses{def:weakly_quasi_complete,def:analytically_irreducible}
  % SOURCE: references/anderson_2014_fulltext.txt, lines 307--339.
  \textit{Source: D.~D. Anderson, Corollary~2,
  references/anderson\_2014\_fulltext.txt, lines 307--339,
  DOI 10.1007/978-1-4939-0925-4\_2.}
  If \(A\) is a one-dimensional Noetherian local domain and the chosen adic
  ideal is the maximal ideal, then \(A\) is weakly quasi-complete if and only
  if its adic completion \(\widehat A\) is a domain.
\end{lemma}

\begin{proof}
  This wrapper directly invokes the explicit published-source boundary
  `anderson_corollary2_part3_external`.  Using the source theorem directly
  keeps the auxiliary generic Hausdorff and Corollary~2(1) boundaries outside
  the final Anderson 8(a) dependency closure.
\end{proof}

\begin{lemma}[Dimension-one analytic criterion]
  \label{lem:dimension_one_weak_criterion}
  \lean{Run202608192034.TODO.dimensionOne_weaklyQuasiComplete_iff}
  \uses{lem:dimension_one_weak_criterion_adic_completion_source}
  % SOURCE: references/anderson_2014_fulltext.txt, lines 307--339.
  \textit{Source: D.~D. Anderson, Corollary~2,
  references/anderson\_2014\_fulltext.txt, lines 307--339,
  DOI 10.1007/978-1-4939-0925-4\_2.}
  This wrapper records the same criterion for a supplied completion target
  \(A^\wedge\).
\end{lemma}

\begin{proof}
  In Lean this declaration is a transport wrapper: once the source supplies
  the relevant analytic-irreducibility criterion for the chosen completion
  target, the theorem returns that equivalence directly.
\end{proof}
